\section{Related Work}
\label{sec:related}

Verifiable ML systems such as SafetyNets~\cite{ghodsi2017safetynets},
zkCNN~\cite{liu2021zkcnn}, vCNN~\cite{lee2024vcnn},
Mystique~\cite{weng2021mystique}, and zkML~\cite{chen2024zkml} all encounter
matrix multiplication as a dominant arithmetic cost.
These systems typically verify matrix products either by arithmetizing the
vector-matrix products directly or by using sumcheck/GKR-style machinery over
layered arithmetic circuits.
$\protocol$ is complementary: it is a dedicated checking layer for the matrix
multiplication component, reducing the checked relation to
$\mathcal{O}(tk+E_{\mathsf{msg}}(k))$ in-circuit work plus
$\mathcal{O}(t\log n+E_{\mathsf{link}}(k))$ auxiliary opening/linking work
under certified encoded commitments.

Dedicated matrix-verification protocols, including
zkMatrix~\cite{cong2024zkmatrix}, DualMatrix~\cite{cong2026dualmatrix},
zkVC~\cite{zhang2025zkvc}, and the coding-theoretic line of Bennett et
al.~\cite{bennett2023matrix}, explore different trade-offs among polynomial
commitments, batching, prover work, sparsity assumptions, and proof size.
These works often package the input-commitment and matrix-verification layers
inside a single committed-matrix system.  Our construction instead exposes a
SNARK-friendly local relation over certified committed codewords and states the
backend requirements explicitly.  This makes the checking layer portable across
commitment backends: systems that already maintain encoded committed inputs can
use $\protocol$ as their online matrix-multiplication checker.
Relative to code-based proof systems such as Ligero~\cite{ames2017ligero},
Brakedown~\cite{golovnev2023brakedown}, Orion~\cite{xie2022orion}, and
FRI~\cite{ben2018fast}, the contribution is not a new general-purpose proof
system; it is the problem-specific reduction that turns the matrix check into
local codeword folds before arithmetization.
The main trade-off is therefore different from both sumcheck/GKR systems and
dedicated matrix-vector-commitment protocols: $\protocol$ verifies multiplication
over certified encoded commitments, targets statement-sound checking rather
than zero knowledge, and delegates codeword binding to the selected backend.

We use Freivalds-in-SNARK as the primary experimental baseline because it
checks the same committed matrix-multiplication statement in the same Groth16
measurement stack and does not require reimplementing a different proof-system
architecture.
Sumcheck/GKR, dedicated matrix VC, and code-based SNARK systems provide
important comparison points, but their costs depend on different input
certification, batching, setup, and transcript models.
We therefore avoid presenting non-normalized running times as direct
head-to-head numbers.
